% !TeX root = ../main.tex
\chapter{Installing EPICS Base and IOCs}
\label{chap:epics_iocs}

\section{Introduction}
A simple IOC was written to interface the Arduino Giga which monitors the six supplies. The following section details how to install EPICS Base and required modules, and compile the IOC. It is written assuming a Linux environment; this particular build was done on Ubuntu 22.04LTS.

\section{EPICS Installation Overview}
	EPICS installations performed by M. Capotosto loosely follow the below directory structure:

	\begin{itemize}
		\item \texttt{EPICS\_BASE=/epics/base} or \texttt{/epics/base\_[version]} or similar
		\item \texttt{SUPPORT=/epics/support} -- or sometimes MODULES instead of SUPPORT is used.
			\par\bigskip
			\quad Modules required for the PSU monitor are:
			\begin{itemize}[topsep=0pt]
				\item pscdrv: Portable Streaming Controller Driver
			\end{itemize}
		\item \texttt{epics/iocs} - each IOC is cloned into this repository
		\item \texttt{epics/systemd-softioc} - This is a utility for managing IOCs as a system service
	\end{itemize}

\section{Installing EPICS Base and Required Modules}
	\begin{enumerate}
		\item Install Dependencies for Build
			\hspace*{2em} Install git, build-essential, and perl:\\ \texttt{sudo apt -y install git build-essential perl libtirpc-dev re2c libevent-dev autoconf automake libtool pkg-config}
		\item Create the EPICS user group, and the empty EPICS directory: \\
			\hspace*{2em} \texttt{sudo groupadd epics}\\
			\hspace*{2em} \texttt{sudo usermod -aG epics \$USER}\\
			\hspace*{2em} \texttt{newgrp epics}\\
			\hspace*{2em} \texttt{sudo mkdir -p /epics}\\
			\hspace*{2em} \texttt{sudo chown \$USER:epics /epics}\\
			\hspace*{2em} \texttt{sudo chmod 2775 /epics}\\
		\item EPICS Base is cloned into \texttt{/epics/base} via:\\
			\hspace*{2em} \texttt{git clone --recursive -b 7.0} \\ \hspace*{4em}\texttt{https://github.com/epics-base/epics-base.git /epics/base}\\
			\hspace*{2em} \texttt{cd /epics/base}\\
		\item Build EPICS:\\
			\hspace*{2em} \texttt{make clean}\\
			\hspace*{2em} \texttt{make}\\

		\item Create the remaining directory structure:\\
			\hspace*{2em} \texttt{mkdir -p /epics/iocs /epics/modules}\\

		\item Clone PSC Driver:
		\begin{enumerate}
			\item \texttt{cd /epics/modules}
			\item \texttt{git clone https://github.com/mdavidsaver/pscdrv.git pscdrv}
			\item \texttt{cd /epics/modules/pscdrv}
			\item Set up required files in \texttt{configure/RELEASE, configure/RELEASE.LOCAL}
			\item Build with \texttt{make clean}, and \texttt{make}
		\end{enumerate}

		\item Clone Manage IOCs Utility:
		\begin{enumerate}
			\item \texttt{cd /epics}
			\item \texttt{git clone https://github.com/NSLS2/systemd-softioc.git}
			\item Navigate to the repo directory \texttt{cd systemd-softioc}
			\item Make the install script executable and run:\\
			\hspace*{2em}\texttt{sudo chmod +x install.sh}\\
			\hspace*{2em}\texttt{sudo ./install.sh}\\
			\hspace*{2em}\texttt{sudo usermod -aG epics softioc}\\
			\hspace*{2em}\texttt{newgrp}\\
			\hspace*{2em}\texttt{sudo chown softioc:epics /epics/iocs}\\
			\hspace*{2em}\texttt{sudo chmod 2775 /epics/iocs}\\
		\end{enumerate}
	\end{enumerate}

\section{Installing the PSU Monitor IOC}
\begin{enumerate}
	\item Change to IOC directory \texttt{cd /epics/iocs}
	\item Clone the repo:\\
	\hspace*{1em} \texttt{git clone https://github.com/capotosto/PS\_Monitor psuMonitorIOC}
	\item \texttt{cd psuMonitorIOC}
	\item Create the \texttt{RELEASE.local} file, \texttt{touch configure/RELEASE.local}
	\item \texttt{echo -e "EPICS\_BASE = /epics/base\textbackslash n}\\
		\hspace*{2em}\texttt{PSCDRV = /epics/modules/pscdrv\textbackslash n}\\
	\item Build the IOC: \texttt{make clean}, \texttt{make}
	\item Update the st.cmd file as needed: \texttt{cd iocBoot/iocpsuMonitor}
	\item Create the st.cmd file used by systemd-softioc: \begin{verbatim}echo -e '#!/bin/bash\n\n
		cd "/epics/iocs/psuMonitorIOC/iocBoot/iocpsuMonitor"\n./st.cmd' > st.cmd\end{verbatim}
	\item Make the new st.cmd executable: \texttt{chmod +x st.cmd}
	\item Find the next available port: \texttt{manage-iocs nextport}
	\item Configure this in the config file: \texttt{nano config}
	\item Set the \texttt{NAME=}, \texttt{PORT=} your next port, \texttt{USER=softioc}, and \texttt{HOST=\$HOSTNAME}
	\item Update the directory path in the root-level st.cmd to the \texttt{iocBoot/iocpsuMonitor/st.cmd}
	\item Enroll in manage-iocs as a service: \texttt{sudo manage-iocs install psuMonitorIOC}
	\item Enable auto-start on boot: \texttt{sudo manage-iocs enable psuMonitorIOC}
	\item Check the status, it should be running! \texttt{manage-iocs status}
	\item And if you want other information on the ioc: \texttt{manage-iocs report}
\end{enumerate}




















